\section{预判大厂的预判：等待、诱惑与转折}

上一章讲资本和商业化，本章看肖弘最重要的战术能力：预判大厂的预判。企业微信 SCRM 的机会来自一个判断：腾讯会治理个人微信外挂，同时把合规需求导向企业微信生态。具体什么时候发生不知道，但可以在成本可承受范围内提前准备。当外挂被打击，准备好的产品就能立刻接住需求。

\begin{figure}[H]
\centering
\includegraphics[width=0.96\textwidth]{predict-giants.png}
\caption{预判大厂的预判：等待、诱惑和大空头式反共识共同构成关键转折。自制概念图，依据 00:42:10--00:52:27 对谈内容整理。}
\end{figure}

\begin{knowledgebox}{读图：确定会发生，不知道何时发生}
图中从观察、预判、等待、诱惑到转折。肖弘的关键方法是：如果你判断某件事大概率会发生，但不知道时间点，就在可承受成本内提前准备。创业难点在于等待期间没有正反馈，还会有定制项目、短期订单和团队动摇。
\end{knowledgebox}

\subsection{大空头式等待：为什么最难}

本节解释等待的痛苦。SCRM 产品上线后，外挂尚未被治理，产品没有用户，团队容易怀疑战略。与此同时，大客户定制、短期收入和其他方向会出现诱惑。肖弘用《大空头》类比：判断正确并不等于马上赚钱，等待期间要承受现金、团队和信心压力。最终当外挂被封，他们当晚写文章、投放传播，迅速成为替代品心智。

\begin{warningbox}{等待期的风险：不是判断错，而是撑不到发生}
创业者常常能看到一个方向，但看不到时间点。方向判断和现金周期、团队耐心、产品完成度必须匹配。太早会消耗，太晚会错过窗口。
\end{warningbox}

\subsection{抵御不相关正反馈}

本节补一个细节。等待期间，大客户定制可能给团队正反馈，销售同学会兴奋，短期现金也诱人。但如果定制项目偏离标准 SaaS 产品，会牵扯研发资源，影响主线迭代。肖弘的做法是拒绝不相关正反馈，把资源留给“等到事件发生时能接住”的产品。这个判断比普通“专注”更具体：不是拒绝所有收入，而是拒绝会破坏战略等待的收入。

\begin{importantbox}{实践经验：不是所有正反馈都该接受}
如果一个反馈来自非战略方向，它可能比负反馈更危险。负反馈会让你警醒，不相关正反馈会让团队误以为找到新路，结果偏离真正等待的窗口。
\end{importantbox}

\subsection{本章小结}

预判大厂的预判，是肖弘第一段创业的核心方法：理解平台会怎样治理生态，提前准备合规替代品，并在等待期抵御短期诱惑。这一方法后来迁移到 AI：预判模型能力什么时候出现，提前把应用容器放在那里。

